
\textbf{Corollaire 7.4.} \emph{Soit $X$ un sch\'ema irr\'eductible de dimension $r\geq2$, qui est projectif sur un corps alg\'ebriquement clos, et soit $H$ un faisceau ample sur $X$. Alors}~:
\begin{enumerate}[label=(\roman*)]
\item \emph{Pour qu'un faisceau inversible $L$ sur $X$ soit num\'eriquement \'equivalent \`a $0$, (il faut et) il suffit que}
\[
   \langle c_1(L)\cdot c_1(H)^{r-1}\rangle=0
   \quad\hbox{\uline{et}}\quad
   \langle c_1(L)^2\cdot c_1(H)^{r-2}\rangle\geq0.
\]
\item \emph{Pour tout faisceau inversible $M$ sur $X$, qui n'est pas num\'eriquement proportionnel \`a $H$, on a}
\[
   \langle c_1(M)^2\cdot c_1(H)^{r-2}\rangle
   <
   \langle c_1(M)\cdot c_1(H)^{r-1}\rangle^2/\langle c_1(H)^r\rangle.
\]
\end{enumerate}

En effet, (ii) r\'esulte tout de suite de (i), d'apr\`es le raisonnement de 7.3. Pour \'etablir (i), soit $Y$ une courbe ferm\'ee, int\`egre, dans $X$, on veut prouver que $\langle c_1(L)\cdot Y\rangle=0$. Plongeons $X$ dans $\mathbb P^N_k$ au moyen de $H^{\otimes n}$ ($n>0$). Soient $I=\bigoplus I_\nu$ l'id\'eal homog\`ene de $Y$, et $\nu_0$ un entier tel que $I_{\nu_0}$ engendre $I_\nu$ pour $\nu\geq\nu_0$. Alors les hypersurfaces $Z$ de degr\'e $\nu_0$ qui contiennent $Y$, n'ont pas de points communs en dehors de $Y$. Raisonnant par r\'ecurrence sur $r=\dim(X)$, on trouve \`a l'aide du th\'eor\`eme de Bertini $r-2$ hypersurfaces $Z_1,\ldots,Z_{r-2}$ de degr\'e $\nu_0$ qui contiennent $Y$, telles que $X'=Z_1\cap\ldots\cap Z_{r-2}\cap X$ soit irr\'eductible. Comme $\langle c_1(L)\cdot Y\rangle=\langle c_1(L_{X'})\cdot Y\rangle$, il suffit de prouver que $L_{X'}$ est num\'eriquement \'equivalent \`a $0$ sur $X'$. Or par la formule de projection, on a que $\langle c_1(L_{X'})\cdot c_1(H_{X'})\rangle=0$ et $\langle c_1(L_{X'})^2\rangle\geq0$ ; donc l'assertion r\'esulte de 7.1.

\textbf{Corollaire 7.5.}
\emph{Soient $k$ un corps alg\'ebriquement clos, $Y$ un $k$-sch\'ema projectif, muni d'un faisceau ample $H$, et $X$ le sch\'ema des z\'eros d'une section $\varphi$ de $H$. On suppose que pour toute composante irr\'eductible $Y^i$ de $Y$, 
